\section{Testing And Verification}
\label{sec:TestingVerification}

The repository now has a meaningful regression net around the scaffolder, the monitoring
path, and the core runtime contracts, even though it is still not a full end-to-end
production validation environment.

\subsection{Generator verification}

Generator verification currently has three layers:
\begin{itemize}
\item dry-run tests that treat scaffold summaries and previews as a contract,
\item generation tests that write selected projects to temporary directories and check the
resulting files and configuration,
\item compile-level tests that run Maven on representative generated projects.
\end{itemize}

This matters because the generator is split across BPMN parsing logic, StringTemplate groups
and multiple helper generators. Simple unit tests are not enough on their own.

\subsection{Representative sample coverage}

The verification set intentionally spans several BPMN areas, including:
\begin{itemize}
\item message and signal events,
\item timers and boundaries,
\item subprocesses and nested subprocesses,
\item call activities,
\item event subprocesses.
\end{itemize}

\subsection{Core runtime contract verification}

Runtime contracts are validated through additional unit test coverage:
\begin{itemize}
\item \texttt{ProcessEvent} serialization round-trips and compact-constructor inference,
\item \texttt{ProcessState} serialization round-trips,
\item \texttt{ScopeCancellationRegistry} cancellation tracking, merging, clearing and null-guard behavior,
\item expanded \texttt{ProcessStateView} projection tests covering cancellation, escalation,
process completion, gateway events, multiple concurrent activities, null-activity edges and
state-key generation.
\end{itemize}

\subsection{Monitoring verification}

The monitoring side is validated through:
\begin{itemize}
\item unit tests around the state-projection model,
\item local startup and packaging checks,
\item demo publishers and scenario runners that can drive the monitoring topology without a
fully generated process runtime.
\end{itemize}

\subsection{CI pipeline}

A GitHub Actions CI workflow runs \texttt{mvn verify} on push and pull requests to the main
branch using JDK 21, ensuring the complete test suite passes on every change.

\subsection{What is not fully covered yet}

The remaining verification gaps are mostly at the integrated runtime level:
\begin{itemize}
\item full end-to-end process executions across all BPMN features,
\item broader load or concurrency-oriented behavior,
\item stronger operational testing of multi-instance monitoring scenarios.
\end{itemize}

For the current project phase, the test strategy is therefore strongest on generator
correctness and local runtime exploration, and lighter on production-style operational proof.
